% concatenated from comeager_hereditary_families_are_big.tex
\documentclass[11pt]{amsart}

\usepackage[T1]{fontenc}
\usepackage[utf8]{inputenc}
\usepackage{lmodern}
\usepackage{microtype}
\usepackage{amsmath,amssymb,amsthm,mathtools}
\usepackage[hidelinks]{hyperref}

\newtheorem{theorem}{Theorem}[section]
\newtheorem{proposition}[theorem]{Proposition}
\newtheorem{lemma}[theorem]{Lemma}
\newtheorem{corollary}[theorem]{Corollary}
\theoremstyle{remark}
\newtheorem{remark}[theorem]{Remark}

\newcommand{\K}{\mathcal K}
\newcommand{\HH}{\mathcal H}
\newcommand{\I}{\mathcal I}
\newcommand{\G}{\mathcal G}
\newcommand{\cl}[1]{\overline{#1}}
\newcommand{\down}[1]{\mathord{\downarrow}\!#1}
\newcommand{\up}[1]{\mathord{\uparrow}\!#1}

\title[Comeager hereditary families]{Comeager hereditary families of compact sets are big}

\author{Peter Banáš}

\address{Faculty of Mathematics and Physics, Charles University, Prague, Czech Republic}
\email{peter.banas149@student.cuni.cz}

\subjclass[2020]{Primary 03E15; Secondary 54B20, 54E52}
\keywords{Vietoris hyperspace, hereditary family, Baire category, Polish space}

\begin{document}

\begin{abstract}
Let $X$ be a Polish space and let $\K(X)$ be its Vietoris hyperspace.
A family $\I\subseteq\K(X)$ is hereditary if it is downward closed under
inclusion. Matheron and Zelený asked whether every comeager hereditary
family in $\K(X)$ contains a dense hereditary $G_\delta$ subfamily.
We give an affirmative answer in ZFC.
\end{abstract}

\maketitle

\section{Introduction}

Let $X$ be a Polish space, and let $\K(X)$ denote the space of all compact
subsets of $X$ equipped with the Vietoris topology.
A family $\I\subseteq\K(X)$ is \emph{hereditary} if
\[
K\subseteq L\in\I \quad\Longrightarrow\quad K\in\I.
\]
Following Matheron and Zelený~\cite{MZ05}, a set
$\mathcal A\subseteq\K(X)$ is called \emph{big} if it contains a dense
hereditary $G_\delta$ subset of $\K(X)$.


Matheron and Zelený asked whether every comeager hereditary family is big
\cite[Problem~2.1]{MZ05}. 

Notice that every comeager family $\I \subseteq\K(X)$ contains a dense $G_\delta$ subset of $\K(X)$. Thus, for a comeager hereditary family, the issue is whether such a subset can be chosen hereditary.

The question was later reiterated in their survey
\cite[Problem~6.9]{MZ07}. They proved the assertion for coanalytic families
\cite[Theorem~2.2]{MZ05}; see also \cite[Theorem~5.11]{MZ07}. Without a
definability assumption, they obtained the conclusion under the additional
set-theoretic hypothesis
\[
\omega_1^{L[x]}<\omega_1
\qquad\text{for every }x\in\omega^\omega
\]
\cite[Theorem~2.9]{MZ05}; see also \cite[Remark~5.12]{MZ07}.
Our main result shows that no additional hypothesis is needed.

\begin{theorem}\label{thm:main}
Let $X$ be a Polish space. Every comeager hereditary family
$\I\subseteq\K(X)$ contains a dense hereditary
$G_\delta$ subset of $\K(X)$.
\end{theorem}

We recall the ideal-theoretic criterion from \cite{MZ05} that will be used to
prove the theorem. For $n\geq1$, put
\[
P(n)=\{K\in\K(X):|K|\leq n\}.
\]
Each $P(n)$ is closed in $\K(X)$. Let $\HH$ be the family of all closed
sets $F\subseteq\K(X)$ such that
\[
F\cap P(n)\text{ is nowhere dense in }P(n)
\qquad\text{for every }n\geq1.
\]
Following \cite[Definition~2.3]{MZ05}, let
\[
\HH^{\mathrm{ext}}
=
\left\{
B\subseteq\K(X):
B\subseteq\bigcup_{m\in\mathbb N}F_m
\text{ for some sequence }(F_m)_{m\in\mathbb N}\subseteq\HH
\right\}.
\]
Thus $\HH^{\mathrm{ext}}$ is the $\sigma$-ideal generated by $\HH$.
The following characterization is \cite[Corollary~2.5]{MZ05}.

\begin{theorem}[Matheron--Zelený]\label{thm:MZcriterion}
For every $\mathcal A\subseteq\K(X)$,
\[
\mathcal A\text{ is big}
\quad\Longleftrightarrow\quad
\K(X)\setminus\mathcal A\in\HH^{\mathrm{ext}}.
\]
\end{theorem}

\section{Preliminary lemmas}

For $A\subseteq\K(X)$, define its hereditary closure by
\[
\down A
=
\{K\in\K(X): (\exists L \in A) (K\subseteq L)\}.
\]

\begin{lemma}\label{lem:open}
If $A\subseteq\K(X)$ is open, then $\down A$ is open.
\end{lemma}

\begin{proof}
Let $K\in\down A$, and choose $L\in A$ with $K\subseteq L$.
If $L=\varnothing$, then $K=\varnothing$, and the isolated point
$\{\varnothing\}$ is contained in $\down A$.

Suppose $L\neq\varnothing$. Choose a basic Vietoris neighbourhood
\[
\langle U_1,\ldots,U_m\rangle
=
\left\{
M\in\K(X):
M\subseteq\bigcup_{i=1}^m U_i
\text{ and }
M\cap U_i\neq\varnothing\text{ for all }i
\right\}
\]
such that
\[
L\in\langle U_1,\ldots,U_m\rangle\subseteq A.
\]
We claim that
\[
\left\{
M\in\K(X):M\subseteq\bigcup_{i=1}^m U_i
\right\}
\subseteq\down A.
\]
Indeed, choose $x_i\in L\cap U_i$ for $1\leq i\leq m$. If
$M\subseteq\bigcup_iU_i$, then
\[
M'=M\cup\{x_1,\ldots,x_m\}
\]
belongs to $\langle U_1,\ldots,U_m\rangle\subseteq A$, while
$M\subseteq M'$. The displayed upper Vietoris set is therefore an open
neighbourhood of $K$ contained in $\down A$.
\end{proof}

Fix a bounded compatible complete metric $d$ on $X$ with $d \leq 1$, and equip the nonempty compact subsets of $X$ with the
associated Hausdorff metric. Extend it to $\K(X)$ by declaring
$d_H(\varnothing,K)=1$ whenever $K\neq\varnothing$. This is a compatible
complete metric for the Vietoris topology.

The next elementary fact is the sequential form of
\cite[Lemma~2.6(i)]{MZ05}.

\begin{lemma}\label{lem:finite-approx}
Let $K,L,L_j\in\K(X)$, where $K\subseteq L$, $K\in P(n)$, and
$L_j\to L$. Then there are $K_j\subseteq L_j$ such that
$K_j\in P(n)$ and $K_j\to K$.
\end{lemma}

\begin{proof}
If $K=\varnothing$, take $K_j=\varnothing$. Otherwise write
$K=\{x_1,\ldots,x_m\}$, where $m\leq n$. Since $L\neq\varnothing$,
the sets $L_j$ are nonempty for all sufficiently large $j$. For such
$j$, choose $x_i^j\in L_j$ so that
\[
d(x_i,x_i^j)\leq d_H(L,L_j)+\frac1j
\qquad (1\leq i\leq m),
\]
and put $K_j=\{x_1^j,\ldots,x_m^j\}$. For the finitely many remaining
indices, put $K_j=\varnothing$. Then $K_j\subseteq L_j$,
$|K_j|\leq n$, and $d_H(K_j,K)\to0$.
\end{proof}

\begin{lemma}\label{lem:closure}
If $A\subseteq\K(X)$ is hereditary, then, for every $n\geq1$,
\begin{equation}\label{eq:closure}
P(n)\cap\cl A=\cl{A\cap P(n)}.
\end{equation}
All closures in \eqref{eq:closure} are taken in $\K(X)$.
\end{lemma}

\begin{proof}
Since $P(n)$ is closed, the inclusion from right to left is immediate.

For the reverse inclusion, let $K\in P(n)\cap\cl A$. If $K=\varnothing$,
then the fact that $\{\varnothing\}$ is open implies
$\varnothing\in A\cap P(n)$.

Otherwise write $K=\{x_1,\ldots,x_m\}$, where $m\leq n$. Let
$\mathcal V$ be an arbitrary basic Vietoris neighbourhood of $K$.
Refining $\mathcal V$, we may choose open sets $V_1,\ldots,V_m$, with
$x_i\in V_i$, such that
\[
K\in\langle V_1,\ldots,V_m\rangle\subseteq\mathcal V.
\]
Since $K\in\cl A$, choose
\[
T\in A\cap\langle V_1,\ldots,V_m\rangle.
\]
For each $i\leq m$, choose $y_i\in T\cap V_i$, and put
$T'=\{y_1,\ldots,y_m\}$. Then $T'\subseteq T$, so heredity gives
$T'\in A$. Moreover, $T'\in P(n)\cap\mathcal V$. Thus every
neighbourhood of $K$ meets $A\cap P(n)$.
\end{proof}


\begin{proposition}\label{prop:local}
Let $O\subseteq\K(X)$ be open and let $F\subseteq\K(X)$ be closed and
nowhere dense. Put
\[
U=\down O,\qquad
V=\down(O\setminus F),\qquad
C=U\setminus V.
\]
Then $\cl C\in\HH$.
\end{proposition}

\begin{proof}
By Lemma~\ref{lem:open}, both $U$ and $V$ are open and hereditary.

We first show that, for every $n\geq1$,
\begin{equation}\label{eq:density}
V\cap P(n)\text{ is dense in }U\cap P(n).
\end{equation}
Fix $K\in U\cap P(n)$, and choose $L\in O$ with $K\subseteq L$.
Since $F$ is nowhere dense and $O$ is open, there are
$L_j\in O\setminus F$ with $L_j\to L$. By
Lemma~\ref{lem:finite-approx}, there are $K_j\subseteq L_j$ such that
$K_j\in P(n)$ and $K_j\to K$. Each $K_j$ belongs to $V$, proving
\eqref{eq:density}.

Suppose towards a contradiction that $\cl C\cap P(n)$ has nonempty
interior in $P(n)$. Choose a nonempty relatively open set
$W\subseteq P(n)$ such that
\[
W\subseteq\cl C.
\]
Since $C\subseteq U$, Lemma~\ref{lem:closure} gives
\[
W\subseteq P(n)\cap\cl U=\cl{U\cap P(n)}.
\]
Since $W$ is nonempty and relatively open in $P(n)$, it follows that
$W\cap U\cap P(n)\neq\varnothing$. This intersection is relatively open
in $U\cap P(n)$, so \eqref{eq:density} yields $W\cap V\neq\varnothing$.

On the other hand, $C\cap V=\varnothing$, and $V$ is open. Hence
\[
\cl C\cap V=\varnothing,
\]
contradicting $W\subseteq\cl C$. Thus $\cl C\cap P(n)$ is nowhere
dense in $P(n)$ for every $n\geq1$, and therefore $\cl C\in\HH$.
\end{proof}

\section{Proof of the main theorem}

\begin{proof}[Proof of Theorem~\ref{thm:main}]
Put
\[
Y=\K(X)
\qquad\text{and}\qquad
B=Y\setminus\I.
\]
Since $\I$ is hereditary, $B$ is upward closed:
\[
K\in B,\quad K\subseteq L
\quad\Longrightarrow\quad
L\in B.
\]
Since $B$ is meager, choose closed nowhere dense sets
$F_r\subseteq Y$, $r\in\mathbb N$, such that
\[
B\subseteq\bigcup_{r\in\mathbb N}F_r.
\]
Let $(O_s)_{s\in\mathbb N}$ be a countable basis for $Y$. For
$r,s\in\mathbb N$, put
\[
U_s=\down O_s,\qquad
V_{r,s}=\down(O_s\setminus F_r),\qquad
C_{r,s}=U_s\setminus V_{r,s}.
\]
By Proposition~\ref{prop:local},
\begin{equation}\label{eq:Hpieces}
\cl{C_{r,s}}\in\HH
\qquad\text{for all }r,s\in\mathbb N.
\end{equation}

We claim that
\begin{equation}\label{eq:cover}
B\subseteq\bigcup_{r,s\in\mathbb N}C_{r,s}.
\end{equation}
Fix $K\in B$, and consider its upper cone
\[
\up K=\{L\in Y:K\subseteq L\}.
\]
The set $\up K$ is closed in $Y$: indeed,
\[
Y\setminus\up K
=
\bigcup_{x\in K}
\{L\in Y:L\subseteq X\setminus\{x\}\},
\]
and the sets on the right are open. Thus $\up K$ is a Baire space.
Since $B$ is upward closed,
\[
\up K\subseteq B\subseteq\bigcup_{r\in\mathbb N}F_r.
\]
By the Baire category theorem, for some $r$, the set
$F_r\cap\up K$ has nonempty relative interior in $\up K$. Hence there
is $s$ such that
\[
\varnothing\neq O_s\cap\up K\subseteq F_r.
\]
The first relation says that some member of $O_s$ contains $K$, so
$K\in U_s$. The second says that no member of
$O_s\setminus F_r$ contains $K$, so $K\notin V_{r,s}$. Therefore
$K\in C_{r,s}$, proving \eqref{eq:cover}.

It follows from \eqref{eq:cover} and \eqref{eq:Hpieces} that
\[
B\subseteq
\bigcup_{r,s\in\mathbb N}\cl{C_{r,s}},
\]
and hence $B\in\HH^{\mathrm{ext}}$. By
Theorem~\ref{thm:MZcriterion}, the family $\I$ is big.
\end{proof}


\section*{Acknowledgements}

I thank Miroslav Zelený for reading the argument and
encouraging its publication.

\begin{thebibliography}{99}

\bibitem[MZ05]{MZ05}
É.~Matheron and M.~Zelený,
\emph{Rudin-like sets and hereditary families of compact sets},
Fund. Math. \textbf{185} (2005), no.~2, 97--116,
\href{https://doi.org/10.4064/fm185-2-1}{doi:10.4064/fm185-2-1}.

\bibitem[MZ07]{MZ07}
É.~Matheron and M.~Zelený,
\emph{Descriptive set theory of families of small sets},
Bull. Symbolic Logic \textbf{13} (2007), no.~4, 482--537,
\href{https://doi.org/10.2178/bsl/1203350880}{doi:10.2178/bsl/1203350880}.

\end{thebibliography}

\end{document}
